
%(BEGIN_QUESTION)
% Copyright 2012, Tony R. Kuphaldt, released under the Creative Commons Attribution License (v 1.0)
% This means you may do almost anything with this work of mine, so long as you give me proper credit

Anta at vi får oppgitt følgende matematiske sannhet:

$$x = y$$

Med andre ord representerer variabelen $x$ nøyaktig samme tallverdi som variabelen $y$. Gitt denne antakelsen må også følgende utsagn være sanne:

$$x + 8 = y + 8$$

$$-x = -y$$

$$9x = 9y$$

$${x \over 25} = {y \over 25}$$

$${1 \over x} = {1 \over y}$$

$$x^2 = y^2$$

$$\sqrt{x} = \sqrt{y}$$

$$e^x = e^y$$

$$\log x = \log y$$

$$x! = y!$$

$${dx \over dt} = {dy \over dt}$$

Forklar det generelle prinsippet som virker her. Hvorfor kan vi endre den grunnleggende likheten $x = y$ på så mange ulike måter, og likevel ende opp med gyldige likheter?

\underbar{file i01306}
%(END_QUESTION)





%(BEGIN_ANSWER)

Prinsippet er dette: Du kan utføre enhver matematisk operasjon på en ligning, så lenge du utfører {\it samme} operasjon på begge sider, på nøyaktig samme måte.

%(END_ANSWER)





%(BEGIN_NOTES)

Dette prinsippet er ikke bare sant, men også svært nyttig når algebraiske uttrykk skal forenkles og omformes.

%INDEX% Mathematics review: basic principles of algebra

%(END_NOTES)
